%
\begin{isabellebody}%
\setisabellecontext{Xor{\isacharunderscore}{\kern0pt}Basics}%
%
\isadelimtheory
%
\endisadelimtheory
%
\isatagtheory
\isakeywordONE{theory}\isamarkupfalse%
\ Xor{\isacharunderscore}{\kern0pt}Basics\isanewline
\ \ \isakeywordTWO{imports}\ Main\isanewline
\isakeywordTWO{begin}%
\endisatagtheory
{\isafoldtheory}%
%
\isadelimtheory
%
\endisadelimtheory
%
\isadelimdocument
%
\endisadelimdocument
%
\isatagdocument
%
\isamarkupsection{Boolean Logic and Stream Cipher Primitives%
}
\isamarkuptrue%
%
\endisatagdocument
{\isafolddocument}%
%
\isadelimdocument
%
\endisadelimdocument
%
\begin{isamarkuptext}%
\subsection{Exclusive OR Formalization}
  Exclusive OR (XOR) serves as a core mathematical primitive for symmetric stream ciphers and One-Time Pads.
  In this module, we formalize XOR logic and verify its algebraic properties.%
\end{isamarkuptext}\isamarkuptrue%
%
\begin{isamarkuptext}%
\paragraph{Definition of XOR}
  We construct XOR from standard conjunction, disjunction, and negation operators: \isa{\isafree{my{\isacharunderscore}{\kern0pt}xor}\ \isafree{a}\ \isafree{b}}.%
\end{isamarkuptext}\isamarkuptrue%
\isakeywordONE{definition}\isamarkupfalse%
\ my{\isacharunderscore}{\kern0pt}xor\ {\isacharcolon}{\kern0pt}{\isacharcolon}{\kern0pt}\ {\isachardoublequoteopen}bool\ {\isasymRightarrow}\ bool\ {\isasymRightarrow}\ bool{\isachardoublequoteclose}\ \isakeywordTWO{where}\isanewline
\ \ {\isachardoublequoteopen}my{\isacharunderscore}{\kern0pt}xor\ a\ b\ {\isasymequiv}\ {\isacharparenleft}{\kern0pt}a\ {\isasymand}\ {\isasymnot}\ b{\isacharparenright}{\kern0pt}\ {\isasymor}\ {\isacharparenleft}{\kern0pt}{\isasymnot}\ a\ {\isasymand}\ b{\isacharparenright}{\kern0pt}{\isachardoublequoteclose}%
\begin{isamarkuptext}%
\paragraph{Evaluation Checks}
  We test basic truth table outcomes using Isabelle's evaluation engine.%
\end{isamarkuptext}\isamarkuptrue%
\isakeywordONE{value}\isamarkupfalse%
\ {\isachardoublequoteopen}my{\isacharunderscore}{\kern0pt}xor\ True\ False{\isachardoublequoteclose}\isanewline
\isakeywordONE{value}\isamarkupfalse%
\ {\isachardoublequoteopen}my{\isacharunderscore}{\kern0pt}xor\ True\ True{\isachardoublequoteclose}%
\begin{isamarkuptext}%
\paragraph{Algebraic Properties}
  First, we establish commutativity for \isa{\isaconst{my{\isacharunderscore}{\kern0pt}xor}}.%
\end{isamarkuptext}\isamarkuptrue%
\isakeywordONE{lemma}\isamarkupfalse%
\ xor{\isacharunderscore}{\kern0pt}comm{\isacharcolon}{\kern0pt}\ {\isachardoublequoteopen}my{\isacharunderscore}{\kern0pt}xor\ x\ y\ {\isacharequal}{\kern0pt}\ my{\isacharunderscore}{\kern0pt}xor\ y\ x{\isachardoublequoteclose}\isanewline
%
\isadelimproof
\ \ %
\endisadelimproof
%
\isatagproof
\isakeywordONE{by}\isamarkupfalse%
\ {\isacharparenleft}{\kern0pt}auto\ simp{\isacharcolon}{\kern0pt}\ my{\isacharunderscore}{\kern0pt}xor{\isacharunderscore}{\kern0pt}def{\isacharparenright}{\kern0pt}%
\endisatagproof
{\isafoldproof}%
%
\isadelimproof
%
\endisadelimproof
%
\begin{isamarkuptext}%
Next, we prove the fundamental property for stream cipher decryption:
  re-applying XOR with key $k$ cancels out the initial encryption step.%
\end{isamarkuptext}\isamarkuptrue%
\isakeywordONE{theorem}\isamarkupfalse%
\ xor{\isacharunderscore}{\kern0pt}cancel{\isacharcolon}{\kern0pt}\ {\isachardoublequoteopen}my{\isacharunderscore}{\kern0pt}xor\ {\isacharparenleft}{\kern0pt}my{\isacharunderscore}{\kern0pt}xor\ m\ k{\isacharparenright}{\kern0pt}\ k\ {\isacharequal}{\kern0pt}\ m{\isachardoublequoteclose}\isanewline
%
\isadelimproof
\ \ %
\endisadelimproof
%
\isatagproof
\isakeywordONE{by}\isamarkupfalse%
\ {\isacharparenleft}{\kern0pt}auto\ simp{\isacharcolon}{\kern0pt}\ my{\isacharunderscore}{\kern0pt}xor{\isacharunderscore}{\kern0pt}def{\isacharparenright}{\kern0pt}%
\endisatagproof
{\isafoldproof}%
%
\isadelimproof
%
\endisadelimproof
%
\begin{isamarkuptext}%
The full cancellation theorem is verified in %
\begin{isabelle}%
\isaconst{my{\isacharunderscore}{\kern0pt}xor}\ {\isacharparenleft}{\kern0pt}\isaconst{my{\isacharunderscore}{\kern0pt}xor}\ \isavar{m}\ \isavar{k}{\isacharparenright}{\kern0pt}\ \isavar{k}\ {\isacharequal}{\kern0pt}\ \isavar{m}%
\end{isabelle}.%
\end{isamarkuptext}\isamarkuptrue%
%
\isadelimtheory
%
\endisadelimtheory
%
\isatagtheory
\isakeywordTWO{end}\isamarkupfalse%
%
\endisatagtheory
{\isafoldtheory}%
%
\isadelimtheory
%
\endisadelimtheory
%
\end{isabellebody}%
\endinput
%:%file=/cygdrive/d/Repositories/She Crypts/formal/week1/theories/Xor_Basics.thy%:%
%:%10=1%:%
%:%11=1%:%
%:%12=2%:%
%:%13=3%:%
%:%27=5%:%
%:%39=8%:%
%:%40=9%:%
%:%41=10%:%
%:%45=14%:%
%:%46=15%:%
%:%48=17%:%
%:%49=17%:%
%:%50=18%:%
%:%52=21%:%
%:%53=22%:%
%:%55=24%:%
%:%56=24%:%
%:%57=25%:%
%:%58=25%:%
%:%60=28%:%
%:%61=29%:%
%:%63=31%:%
%:%64=31%:%
%:%67=32%:%
%:%71=32%:%
%:%72=32%:%
%:%81=35%:%
%:%82=36%:%
%:%84=38%:%
%:%85=38%:%
%:%88=39%:%
%:%92=39%:%
%:%93=39%:%
%:%102=42%:%
%:%105=42%:%
%:%113=44%:%
